\section{Analytic construction of $\mathcal{P}$}

\paragraph*{Notation and disambiguation.}
The symbol $\ii=\mathrm{i}$ denotes the imaginary unit. Integer indices are consistently written as $i$; the context will always disambiguate from the imaginary unit. For lattices generated by an index, the explicit set notation $\{\, i k : k\in\mathbb{Z}\,\}$ is used instead of $i\mathbb{Z}$ to avoid ambiguity. Asymptotic symbols $O(\cdot)$ and $o(\cdot)$ have their standard meaning; unless stated otherwise, bounds are uniform on compact $x$–intervals $K\subset(0,\infty)$ for fixed integer parameters.

\subsection{Motivation}\label{subsec:motivation}
The construction uses an analytic representation of trial division. For a real number $x$ and an integer $i \ge 2$, divisibility can be expressed by the quotient
\[
Q(x,i) = \frac{\sin^2(\pi x)}{\sin^2(\pi x/i)}.
\]
The numerator vanishes if and only if $x$ is an integer. If $x$ is an integer, the denominator vanishes if and only if $i$ divides $x$. The quotient thus becomes indeterminate of the form $0/0$ precisely at integer arguments $x=n$ where $i$ is a divisor of $n$.

The indeterminacy is removed by the corresponding Fejér cosine polynomial $F(x,i)$. By Fejér’s identity the polynomial equals the quotient away from the removable singularities, which provides a well-defined extension at those points. For integers $n$, one has $F(n,i)=i^2$ if $i\mid n$ and $F(n,i)=0$ otherwise.

\subsection{Regularisation}
By the Fejér kernel identity~\cite{zygmund2002}
\[
\frac{\sin^{2}(r y)}{\sin^{2}y}=r+2\sum_{k=1}^{r-1}(r-k)\cos(2k y)\qquad (r\in\mathbb N),
\]
and with \(r=i\), \(y=\pi x/i\), one has the pointwise identity
\[
F(x,i)=i+2\sum_{k=1}^{i-1}(i-k)\cos\!\left(\frac{2\pi k x}{i}\right)
=\left(\frac{\sin(\pi x)}{\sin(\pi x/i)}\right)^{\!2},
\]
which is understood as the holomorphic continuation across the removable singularities. This yields a finite cosine polynomial that matches the quotient for \(x\notin\mathbb{Z}\) and reproduces the integer divisibility pattern.

\subsection{Definition}
The regularized function is constructed by summing the Fejér-kernel-based terms over all potential divisors up to $\sqrt{x}$.
Define $F(x,i)$ for $i\in\mathbb{N}$, $i\ge 2$, by
\begin{equation}\label{eq:Fxi}
F(x,i) = i+2\sum_{k=1}^{i-1}(i-k)\cos\left(\frac{2\pi k x}{i}\right).
\end{equation}

% remark: Global nonnegativity is stated in Proposition~\ref{prop:fejer-properties}(c).

\begin{definition}\label{def:Px}
For $x>1$ the function $\mathcal{P}(x)$ is defined as
\begin{equation}\label{eq:Px}
\mathcal{P}(x)=\frac{1}{x}\sum_{i=2}^{\lceil\sqrt{x}\rceil} F(x,i),
\end{equation}
and $\mathcal{P}(x)=0$ is set for $x\le1$.
\end{definition}

\begin{remark}[On the choice of summation limit]
The ceiling function $\lceil\sqrt{x}\rceil$ is used instead of the more common floor function to ensure $C^1$-smoothness. As shown in Proposition~\ref{prop:C1}, at integer squares $x=m^2$, the new term $F(x, m+1)$ and its first derivative vanish, which guarantees a smooth transition. The floor function would be worse: at $x=m^2$ the new index $i=m$ enters with $F(m^2,m)=m^2$, so that the resulting function would already jump in value, by exactly $1$, at every square.
\end{remark}

Write $N(x)=\lceil\sqrt{x}\rceil$ for the upper summation limit; then \eqref{eq:Px} is $\mathcal{P}(x)=\frac{1}{x}\sum_{i=2}^{N(x)} F(x,i)$ for $x>1$.

\begin{remark}[Integer evaluation]\label{rem:integer-eval}
For an integer $n\ge 2$, Proposition~\ref{prop:fejer-properties} implies that $F(n,i)=i^2$ if $i\mid n$ and $F(n,i)=0$ otherwise. The definition in \eqref{eq:Px} therefore yields the exact evaluation
\[
\mathcal P(n)=\frac{1}{n}\sum_{\substack{d\mid n \\ 2\le d\le \lceil\sqrt{n}\rceil}} d^2.
\]
This sum vanishes exactly when $n$ has no divisors in the summation range. For $n>2$ this characterizes primality, but a direct evaluation requires $\Theta(\sqrt n)$ work and thus entails no algorithmic advantage over elementary trial division. At integer squares $n=m^2$, the limit $\lceil\sqrt{n}\rceil=m$ coincides with $\lfloor\sqrt{n}\rfloor$.
\end{remark}

\begin{proposition}[Properties of the Fejér term]\label{prop:fejer-properties}
For an integer $i \ge 2$, the function $F(\cdot,i)$ defined in \eqref{eq:Fxi} is an entire function with the following properties:
\begin{enumerate}
    \item[(a)] \textbf{Closed form and entire extension.} The trigonometric polynomial $F(\cdot,i)$ is entire. The map
    \[
    z\ \longmapsto\ \left(\frac{\sin(\pi z)}{\sin(\pi z/i)}\right)^{\!2}
    \]
    is meromorphic on $\mathbb{C}$ with removable singularities at the lattice points $z=ik$ ($k\in\mathbb{Z}$), because numerator and denominator vanish to the same order there. On the punctured plane $\mathbb{C}\setminus\{ik:k\in\mathbb{Z}\}$ the Fejér identity (with $y=\pi z/i$) yields pointwise equality with the cosine polynomial $F(\cdot,i)$. Removing the removable singularities produces the unique entire extension of this meromorphic function; by the identity theorem the resulting entire function coincides identically with $F(\cdot,i)$ on $\mathbb{C}$.
    \item[(b)] \textbf{Integer values.} For any $n \in \mathbb{Z}$,
    \[
    F(n,i) \;=\; \begin{cases} i^2, & \text{if } i \mid n, \\ 0, & \text{if } i \nmid n. \end{cases}
    \]
    \item[(c)] \textbf{Positivity on $\mathbb{R}$.} For all $x\in\mathbb{R}$, $F(x,i) \ge 0$, with strict inequality for $x\notin\mathbb{Z}$.
\end{enumerate}
\end{proposition}

\begin{proof}
The trigonometric polynomial in \eqref{eq:Fxi} defines an entire function. For $z \notin \{\,i k: k\in\mathbb{Z}\,\}$, Fejér's identity gives
\[
F(z,i)=\left(\frac{\sin(\pi z)}{\sin(\pi z/i)}\right)^{\!2}.
\]
The right-hand side is holomorphic on $\mathbb{C}\setminus\{\,i k: k\in\mathbb{Z}\,\}$ and has removable singularities at the points $z=i k$. Removing these poles yields an entire function that agrees with the trigonometric polynomial on a nonempty open set; the identity theorem then gives global equality.

For (b), if $i \nmid n$ then $\sin(\pi n/i)\ne 0$ while $\sin(\pi n)=0$, hence $F(n,i)=0$. If $i \mid n$, write $n=i m$ and apply L’Hôpital's rule to obtain $F(n,i)=i^2$.

For (c), if $x\notin\mathbb{Z}$ then $\sin(\pi x/i)\ne 0$ for every $i\ge 2$, so $F(x,i)=\bigl(\sin(\pi x)/\sin(\pi x/i)\bigr)^2>0$. At $x\in\{i k\}$ the quotient has a removable singularity and the extension yields the stated nonnegativity, with $F(x,i)=i^2>0$ when $x=i k$.
\end{proof}

\begin{proposition}[Resonant partial-fraction identity and convergence]\label{prop:RPF}
Let $i\in\mathbb{N}$, $i\ge 2$, and $z\in\mathbb{C}$. On the punctured plane $\mathbb{C}\setminus\{\,i k: k\in\mathbb{Z}\,\}$ one has
\[
\left(\frac{\sin(\pi z)}{\sin(\pi z/i)}\right)^{\!2}
=\frac{i^{2}}{\pi^{2}}\,\sin^{2}(\pi z)\,\sum_{k\in\mathbb{Z}}\frac{1}{(z-i k)^{2}},
\]
and the series on the right converges locally uniformly on compact subsets of this domain, so the displayed equality holds pointwise on $\mathbb{C}\setminus\{i k:k\in\mathbb{Z}\}$. Both sides therefore represent the same meromorphic function with removable singularities at $z=i k$. Removing these singularities produces the unique entire extension; by the identity theorem this entire extension coincides identically with the Fejér cosine polynomial $F(z,i)$.
\end{proposition}

\begin{proof}
The classical expansion $\pi^{2}\csc^{2}(\pi u)=\sum_{k\in\mathbb{Z}}(u-k)^{-2}$ converges locally uniformly on compact subsets of $\mathbb{C}\setminus\mathbb{Z}$ (Whittaker–Watson~\cite{whittaker1927}, Chapter VII, §2, cosecant expansion). With $u=z/i$,
\[
\csc^{2}\!\Bigl(\frac{\pi z}{i}\Bigr)=\frac{i^{2}}{\pi^{2}}\sum_{k\in\mathbb{Z}}\frac{1}{(z-i k)^{2}}
\]
on $\mathbb{C}\setminus\{\,i k: k\in\mathbb{Z}\,\}$. For any compact $K\subset\mathbb{C}\setminus\{\,i k: k\in\mathbb{Z}\,\}$ there exist constants $C_K,c_K>0$ such that $|z-i k|\ge c_K\,|k|$ for all $z\in K$ and $|k|$ sufficiently large. Hence
\[
\sum_{|k|>K_0}\frac{1}{|z-i k|^{2}}\ \le\ \frac{C_K}{c_K^{2}}\sum_{|k|>K_0}\frac{1}{k^{2}},
\]
and the Weierstrass M-test yields uniform convergence on $K$. Multiplication by $\sin^{2}(\pi z)$ preserves local uniform convergence on compacta disjoint from $\{i k\}$ because $\sin^{2}(\pi z)$ is holomorphic and bounded on each such compact set. This gives the identity
\[
\left(\frac{\sin(\pi z)}{\sin(\pi z/i)}\right)^{\!2}
=\frac{i^{2}}{\pi^{2}}\,\sin^{2}(\pi z)\,\sum_{k\in\mathbb{Z}}\frac{1}{(z-i k)^{2}}
\]
on $\mathbb{C}\setminus\{\,i k: k\in\mathbb{Z}\,\}$. At the points $z=i k$ both sides have removable singularities (the quotient has a $0/0$ form and the series side is a bounded holomorphic factor times a square-summable tail). Removing these singularities yields the unique entire extension of the common meromorphic function. Since both sides agree with the trigonometric polynomial $F(z,i)$ on a nonempty open subset of $\mathbb{C}$, the identity theorem implies global equality with $F(z,i)$.
\end{proof}


\begin{definition}[Nearest-lattice notation]\label{def:nearest-lattice}
For $x\in\mathbb{R}$ and an integer index $i\ge2$, set $m:=\mathrm{round}(x/i)$ with the fixed convention $\mathrm{round}(u):=\lfloor u+\tfrac12\rfloor$ (ties to the larger integer), and write $t:=x-im$. Then $|t|\le i/2$. In particular, $\sin(\pi x)=\pm \pi t+O(t^{3})$ and hence
\[
F(x,i)=\left(\frac{\sin(\pi x)}{\sin(\pi x/i)}\right)^{\!2}
=i^{2}+O\!\bigl(i^{2}(\pi t)^{2}\bigr)\qquad(t\to0),
\]
which recovers $F(im,i)=i^{2}$ by the removable-singularity extension.
\end{definition}

\begin{lemma}[Nearest-pole dominance bounds]\label{lem:nearest-pole-bounds}
Fix $i\ge2$ and $x\in\mathbb{R}$. Set $m:=\mathrm{round}(x/i)$ and $t:=x-i m$. The Fej\'er term can be written as $F(x,i) = \frac{i^2}{\pi^2}\sin^2(\pi x)\,\Sigma(x,i)$, where the pole sum $\Sigma(x,i) = \sum_{k\in\mathbb{Z}}(x-i k)^{-2}$ satisfies
\[
\frac{1}{t^2} \;\le\; \Sigma(x,i) \;\le\; \frac{1}{t^2} + \frac{\pi^2}{i^2}.
\]
\end{lemma}

\begin{proof}
The representation follows from Proposition~\ref{prop:RPF}. The sum $\Sigma(x,i)$ can be written as
\[
\Sigma(x,i) = \frac{1}{t^2}+\sum_{r=1}^{\infty}\left[\frac{1}{(x-i(m+r))^{2}}+\frac{1}{(x-i(m-r))^{2}}\right].
\]
The lower bound is obtained by discarding the non-negative tail sum. For the upper bound, $|t|\le i/2$ by definition, so
\[
|x-i(m\pm r)| = |t\mp ri| \ge ri - |t| \ge (r-\tfrac12)i\qquad(r\ge1).
\]
This gives the tail bound
\[
\sum_{r=1}^{\infty}\left[\dots\right] \le \frac{2}{i^2}\sum_{r=1}^{\infty}\frac{1}{(r-\tfrac12)^2} = \frac{2}{i^2}\cdot\frac{\pi^2}{2} = \frac{\pi^2}{i^2},
\]
where the sum identity follows from the standard cosecant series evaluated at $z=1/2$.
\end{proof}

\begin{remark}[Integer arguments and removable singularities]
If $x=n\in\mathbb{Z}$, the inequalities in Lemma~\ref{lem:nearest-pole-bounds} are to be understood in the removable-singularity sense. If $i\nmid n$, then $t\neq0$ and the bounds hold literally (while $F(n,i)=0$ because $\sin(\pi n)=0$). If $i\mid n$, write $t=x-n$; since $\sin(\pi x)\sim \pi(x-n)$ as $x\to n$, one has
\[
\frac{i^{2}}{\pi^{2}}\frac{\sin^{2}(\pi x)}{t^{2}}\ \longrightarrow\ i^{2}\ =\ F(n,i),
\]
which shows that both bounds converge to the common value at $x=n$.
\end{remark}

\begin{remark}[An identity on the divisor lattice]
For integer arguments $n$, the normalized Fejér terms satisfy a multiplicative property on the divisor lattice. Since $F(n,r)/r^2$ is the indicator function $\mathbf{1}_{r\mid n}$ for integers, the identity $\mathbf{1}_{i\mid n}\mathbf{1}_{j\mid n}=\mathbf{1}_{\operatorname{lcm}(i,j)\mid n}$ implies
\[
\frac{F(n,i)}{i^2}\cdot \frac{F(n,j)}{j^2}\;=\;\frac{F\bigl(n,\operatorname{lcm}(i,j)\bigr)}{\operatorname{lcm}(i,j)^2}
\]
for any integers $i,j\ge 2$. This property is specific to integer inputs.
\end{remark}

\begin{remark}[Classical references]
Standard facts on Fejér kernels, cosine–polynomial representations, and Cesàro means are classical; see Zygmund~\cite[Chap.~I, §5; Chap.~III]{zygmund2002}, Katznelson~\cite[Chap.~I]{katznelson2004}, and Stein–Shakarchi~\cite[Chap.~2]{stein2003}. Nonnegativity follows from the Fejér–Riesz factorization for nonnegative trigonometric polynomials; cf.\ \cite[Chap.~V]{zygmund2002} and \cite[Chap.~II]{katznelson2004}.
\end{remark}

% ---------------------------------------------------------------------------

